\section{Uniform arithmetic estimates}
\label{p1:arith:appendix}

This section supplies the uniformity needed when the window length grows.
The degree parameters in the determinant arguments are fixed before $N$
tends to infinity.  Dependence on such fixed parameters is permitted;
dependence on a squarefree coefficient of size $N^{O(1)}$ is not.

\subsection{Split products, defects, and the regular set}

Products in short intervals have a long history; see
Erd\H{o}s and Turk~\cite{erdosturk1984}.  The polynomial square-root
argument for a fixed even number of distinct factors appears in the
proof of Munsch, Shparlinski, Sun, and Xiao~\cite[Lemma~3.3]{munsch2026}.
We record an explicit bound uniform in the degree, as required when the
number of selected factors grows.  No constant depending implicitly on
that degree is used in the following lemma.

\begin{lemma}[A split-product bound uniform in the degree]\label{p1:arith:runge}
Let $d=2k\geq2$ and let $\gamma_1,\ldots,\gamma_d$ be distinct integers
with $|\gamma_i|\leq R$, $R\geq1$.  If $U\geq2R$ is an integer and
$\prod_i(U+\gamma_i)$ is a square, then
\[
 U\leq(64R)^d.
\]
\end{lemma}

\begin{proof}
Write $G(T)=\prod_i(T+\gamma_i)$ and
\[
 \phi(z)=\prod_i(1+\gamma_i z)^{1/2}=\sum_{j\geq0}c_jz^j,
 \qquad P(T)=\sum_{j=0}^k c_jT^{k-j}.
\]
The branches equal one at the origin.  Cauchy's bound on $|z|=(2R)^{-1}$
gives $|c_j|\leq(3/2)^{d/2}(2R)^j$ and, for $j\leq k$,
$|c_j|\leq2^dR^j$.  Comparing coefficients in $\phi^2$ inductively gives
$2^{2j-1}c_j\in\mathbb Z$ for $j\geq1$, so $2^dP\in\mathbb Z[T]$.
For $U>4R$ the tail satisfies
\[
 |\sqrt{G(U)}-P(U)|\leq4\,3^{d/2}R^{k+1}/U.
\]
If this difference is nonzero, it is at least $2^{-d}$, whence
$U\leq(8R)^d$.  If it vanishes, the nonzero integer polynomial
$Q=2^{2d}(P^2-G)$ has root $U$.  It is nonzero because $G$ has simple roots.
Keeping the homogeneous powers of $R$ in the coefficient products gives
\[
 |[T^{d-j}](P^2-G)|\leq(d+2)4^dR^j\leq8^dR^j,
 \qquad H(Q)\leq(32R)^d.
\]
The integral-root bound gives $U\leq1+H(Q)\leq(64R)^d$.
The range $2R\leq U\leq4R$ already satisfies the conclusion.
\end{proof}

The following coefficient-uniform count is classical; see
Cilleruelo and Garaev~\cite[Section~3, Lemma~4 and Proposition~1]{cilleruelogaraev2011}
and the directly comparable formulation of
Tao~\cite[Lemma~2.10]{tao2026anatomy}.  When $ab$ is not a square,
the substitution $(au)^2-abv^2=ah$ reduces it to the former lemma;
the square case is a divisor count.  We include a proof to keep the
polynomial coefficient range explicit.

\begin{lemma}[Uniform quadratic equation count]\label{p1:arith:pell}
Fix $A_0>0$.  Uniformly for positive integers $a,b$ and nonzero integers $h$
with $a,b,|h|\leq N^{A_0}$, the number of positive solutions to
\[
 au^2-bv^2=h,\qquad u,v\leq N^{A_0},
\]
is $N^{o_{A_0}(1)}$.
\end{lemma}

\begin{proof}
If $ab$ is a square, multiply the equation by $a$ and factor the resulting
difference of squares; the divisor bound suffices.  Otherwise write
$ab=Ds^2$ with $D$ squarefree and nonsquare.  Each solution gives the
algebraic integer
$\xi=au+sv\sqrt D$ of norm $ah$ in the maximal order of
$\mathbb Q(\sqrt D)$.  Its principal ideal has norm $|ah|$.
The number of ideals of this norm is at most $\tau(|ah|)$: at a split prime
there are at most $e+1$ ideals of norm $\ell^e$, and at an inert or ramified
prime there are no more.  Elements generating the same ideal differ by a
unit.  A real quadratic unit greater than one is at least
$(1+\sqrt5)/2$, as follows from its integral trace and norm $\pm1$.
The positive values of $\xi$ arising here lie in $[1,N^{O(A_0)}]$;
there are therefore $O_{A_0}(\log N)$ possibilities in each ideal class
of generators.  Finally $\tau(|ah|)=N^{o_{A_0}(1)}$.  This argument makes
the coefficient uniformity explicit and requires no bound for a fundamental
unit or a class number.
\end{proof}

\begin{lemma}[Local and global defect counts]\label{p1:arith:defects}
Suppose $B\asymp\log N$ and $H\asymp B$, with fixed comparison constants.
Any interval of diameter $O(B)$ in $[cN,CN]$ contains
$O(B/\log B)$ $H$-defective integers.  Moreover
\begin{align}
 \#\{n\leq CN:K_H(n)=1\}&\leq N^{1/2+o(1)},\label{p1:arith:onedfect}\\
 \#\{x\in I_N:V_x\text{ contains two }H\text{-defective integers}\}
   &\leq N^{o(1)}.\label{p1:arith:twodefect}
\end{align}
All statements remain valid after fixed dilations and translations of the
intervals.
\end{lemma}

\begin{proof}
For $D$ defective integers in a short interval, form the binary matrix with
their small-prime valuation columns and an additional parity row.  Its
kernel has redundancy at most $\pi(H)+1$.  A nonzero word selects an even
square product, so Lemma~\ref{p1:arith:runge} gives minimum distance
$\delta\gg\log N/\log B$.  Choose
$t\asymp\log N/\log B$ with $2t<\delta$.  The Hamming bound gives
\[
 \sum_{j=0}^t\binom Dj\leq2^{\pi(H)+1}.
\]
If $D\geq2t$, use $\binom Dt\geq(D/t)^t$.  Since
$(\pi(H)+1)/t=O(1)$, this proves $D=O(B/\log B)$; the other case is
immediate.

Write a defective integer uniquely as $n=a^2c$, where $c$ is squarefree
and $H$-smooth.  Then
\[
 \#\{n\leq CN:K_H(n)=1\}
 \leq\sqrt{CN}\prod_{\ell\leq H}(1+\ell^{-1/2})
 \leq\sqrt{CN}\exp(O(\sqrt H)).
\]
For two distinct defective values in one window, write
$x+i=c_1a^2$, $x+j=c_2b^2$.  The offsets and the squarefree smooth parts
have $B^2 4^{\pi(H)}=N^{o(1)}$ choices.  All their heights are polynomial
in $N$, and $c_1a^2-c_2b^2=i-j\ne0$.
Lemma~\ref{p1:arith:pell} supplies $N^{o(1)}$ possibilities.  This proves
\eqref{p1:arith:twodefect} and the stated uniform versions.
\end{proof}

\subsection{Two small kernels with their smooth coefficients retained}

The following form of the approximate determinant method is needed because
the small-prime parts are fixed separately and may be polynomially large.
It is an extension of the proof of \cite[Sections~2.1--2.3]{reuss2014},
rather than an application of its fixed-coefficient statement.

\begin{lemma}[Uniform determinant estimate]\label{p1:arith:reuss}
Fix $A_0,\delta,\varepsilon>0$.  Suppose $1\leq c_1,c_2\leq3N$,
$0<|h|\leq(\log N)^{A_0}$, and
\[
 c_1D^2U\asymp c_2E^2V\asymp N,\quad
 D,E\ll\sqrt N,\quad c_1c_2(DE)^2\geq N,
 \quad\log(DE),\log(UV)\geq\delta\log N.
\]
The comparison constants are fixed.  The number of positive solutions
$d\asymp D,e\asymp E,u\asymp U,v\asymp V$ to
$c_2e^2v-c_1d^2u=h$ is
\begin{equation}\label{p1:arith:reussbound}
 \ll_{A_0,\delta,\varepsilon}N^\varepsilon
 \min\{(DE\mathcal M)^{1/2}+D+E,
          (UV\mathcal M)^{1/2}+U+V\},
 \qquad
 \log\mathcal M=\frac98\frac{\log(DE)\log(UV)}{\log N}.
\end{equation}
Fixed or polylogarithmic factors in the lower bound for $c_1c_2(DE)^2$
can be absorbed by decreasing the error margin.
\end{lemma}

\paragraph{Parameters in the determinant argument.}
Here and in the proof, height of a polynomial means the maximum absolute
value of its integer coefficients.  The ledger below distinguishes the
fixed auxiliary degrees from the coefficients that vary with $N$.
The symbols $A,B_0,\kappa,L_1,T_1$ are specified in Steps~2--3.
\begin{center}
\small
\begin{tabularx}{\textwidth}{@{}>{\raggedright\arraybackslash}p{.23\textwidth}>{\raggedright\arraybackslash}p{.19\textwidth}>{\raggedright\arraybackslash}p{.22\textwidth}>{\raggedright\arraybackslash}X@{}}
\toprule
Object and dimension & Degree & Point or box scale & Coefficient height\\
\midrule
Input hypersurface in $\mathbb A^4$, dimension $3$
& Total degree $3$
& $(D,E,U,V)$, each $O(N)$
& $c_1,c_2\leq3N$; $|h|\leq(\log N)^{A_0}$\\
Auxiliary curve in $\mathbb P^1\!\times\!\mathbb P^1$, dimension $1$
& Bidegree $(a,b)\leq(A,B_0)$, fixed
& Primitive pairs of heights $L_1,T_1\leq N^C$
& $\|G\|\leq N^\kappa$ after the lattice changes\\
Eliminant curve in $\mathbb A^2$, when $a=1$ or $b=1$
& Bounded in terms of $A,B_0$
& The same primitive boxes
& Resultants and nonzero right-hand sides $N^{O(\kappa)}$\\
\bottomrule
\end{tabularx}
\end{center}
All displayed height exponents, including $C$, are fixed, uniformly over
the input boxes.
The relative chart has errors $\alpha=O(M^{-1})$ and
$\beta=O(|h|/N)$, with $M=N^\zeta\mathcal M$; no coefficient ratio
is charged to the latter error.  Choose the final $\varepsilon$ first,
then a sufficiently small $\zeta$, then $A,B_0$ and $\kappa$, and finally
the curve-count exponent $\epsilon'$ so that
$(C+\kappa)\epsilon'\leq\zeta$.  Only after these choices does $N$
tend to infinity.  In particular, the factor $\|G\|^{\epsilon'}$
in Heath-Brown's bound is explicitly included in this budget.
The coefficient uniformity in this ledger is established below;
Reuss's stated theorem allows its constant to depend on the fixed shift
$h$ and does not itself assert this extension.

\begin{proof}
The proof has four steps.  They identify separately the determinant
scale, the uniform count on each auxiliary curve, and the two lattice
summations.  Only the last step sums over the short intervals.

\medskip\noindent
\textbf{Step 1: primitive pairs and relative normalization.}
If $g=\gcd(d,e)$ and $q=\gcd(u,v)$, then $g^2q\mid h$.
Divide by $g^2q$ and replace the four dyadic scales accordingly.
Both pairs are now primitive.  The effective size $N/(g^2q)$ and the
hypotheses on the scales change only by polylogarithmic factors;
in particular the lower bound for $c_1c_2(DE)^2$ has the permitted
polylogarithmic loss.  There are only polylogarithmically many reductions,
which will be summed at the end.  We keep the notation $N,D,E,U,V,h$
for a reduced box.

Set $s=d/e$, $t=v/u$.  The exact curve approximation is
\[
 t=(c_1/c_2)s^2+h/(c_2e^2u).
\]
On a relative interval of length $M^{-1}$ write
\[
 s=s_0(1+\alpha),\qquad
 t=t_0(1+2\alpha+\alpha^2+\beta),\qquad
 t_0=(c_1/c_2)s_0^2\asymp V/U.
\]
Then $|\alpha|\ll M^{-1}$ and $|\beta|\ll|h|/N$, uniformly in
$c_1,c_2$.  The relevant vertical error is relative to $t_0$;
its bound therefore contains no additional factor $c_1/c_2$.

\medskip\noindent
\textbf{Step 2: the determinant with variable coefficients.}
For the auxiliary monomials $s^at^b$, $0\leq a\leq A$,
$0\leq b\leq B_0$, the row denominator is $e^Au^{B_0}$.
Extracting $s_0^at_0^b$ from the columns leaves the same logarithmic
denominator cost as in the normalized calculation:
\[
 A\log E+B_0\log U+\frac A2\log(D/E)+\frac{B_0}2\log(V/U)
 =\frac A2\log(DE)+\frac{B_0}2\log(UV).
\]
Take $M=N^\zeta\mathcal M$, with a small fixed $\zeta>0$, and use
$N/|h|$ for the vertical approximation scale.  The strict determinant
inequality in \cite[Section~2.1, Lemma~7]{reuss2014} then holds, with
degrees fixed in terms of $\zeta,\delta$.  The resulting integer
polynomial has coefficient height $N^{\kappa(\zeta,\delta)}$.
The exponent is uniform in $c_1,c_2$: their sizes, and the rational
column scalings, are polynomial in $N$, while the degrees are fixed.

To check the allowed scale, put $z=\log(DE)/\log N$ and
$w=1-\log(c_1c_2)/(2\log N)$.  Then
\[
 \frac{\log(UV)}{\log N}=2w-2z+O(1/\log N),\qquad
 \frac{\log\mathcal M}{\log N}=\frac94z(w-z)+o(1)
                    \leq\frac9{16}+o(1).
\]
Thus $M\leq N/|h|$ for sufficiently small $\zeta$.
The two logarithmic size ratios stay in a compact subset of the positive
quadrant, so the degree choices are uniform over all boxes.
The parameter order is: choose $\zeta$ from the final error margin,
then the auxiliary degrees and the height exponent $\kappa$, and only
then the internal exponent in the curve count below.

\medskip\noindent
\textbf{Step 3: fibres and the count on one auxiliary curve.}
Apply the two unimodular lattice changes of
\cite[Section~2.2]{reuss2014} to $(d,e)$ and $(v,u)$.
Write the resulting primitive pairs as $\lambda=(\lambda_1,\lambda_2)$
and $\tau=(\tau_1,\tau_2)$, in boxes of side lengths
$L_1\geq L_2$ and $T_1\geq T_2$, respectively.  The changes of
coordinates have polynomial height; enlarging $\kappa$ absorbs their
contribution to the heights of the original and auxiliary equations.
The original equation still has bidegree $(2,1)$.

An auxiliary factor of degree zero in $(d,e)$ fixes only boundedly
many primitive $(u,v)$; Lemma~\ref{p1:arith:pell} counts the remaining
quadratic equation uniformly.  For degree zero in $(u,v)$, fix $(d,e)$
and set $g_f=\gcd(c_1d^2,c_2e^2)$.  Solvability requires $g_f\mid h$.
The step of the resulting one-dimensional lattice in $u$ is
$c_2e^2/g_f$; the step in $v$ is $c_1d^2/g_f$.  Its intersection
with the box has at most
\[
 O\left(1+\frac{Ng_f}{c_1c_2(DE)^2}\right)
 \leq O\left(1+\frac{N|h|}{c_1c_2(DE)^2}\right)
\]
points.  This is polylogarithmic by the stated lower bound.

For an absolutely irreducible factor of bidegree $(a,b)$,
$a,b\geq1$, \cite[Lemma~2]{heathbrown2012} bounds primitive solutions by
$O_{a,b,\epsilon'}(X^{2/b+\epsilon'}\|G\|^{\epsilon'})$ when the
first pair has height at most $X$; the second pair has no height
restriction.  Thus $b\geq2$ gives the required bound in $L_1$, and
$a\geq2$ gives the bound in $T_1$ after exchanging the pairs.
Choose $\epsilon'$ after $\kappa$ so that the coefficient-height
factor costs at most the chosen power $N^{O(\zeta)}$.

The missing bound when one bidegree is one follows directly by
elimination.  In bidegree $(1,b)$ write
$\lambda_1F_1(\tau)+\lambda_2F_2(\tau)=0$.
Primitivity gives
$q\lambda_1=F_2(\tau)$, $q\lambda_2=-F_1(\tau)$ for a nonzero integer
$q$.  The two dehomogenizations of the coprime forms $F_1,F_2$ give
nonzero resultants $R,R'$ of height $N^{O(\kappa)}$ and
\[
 q\mid R\tau_1^m,\qquad q\mid R'\tau_2^{m'}.
\]
Since $(\tau_1,\tau_2)=1$, prime by prime $q$ divides $RR'$.
There are $N^{o(1)}$ possibilities, uniformly in the coefficients.
Substitution into the original equation gives
$F_3(\tau_1,\tau_2)=hq^2$, where $F_3$ is homogeneous of fixed degree
$2b+1$.  If $F_3$ is zero there is no solution.  Otherwise
$F_3-hq^2$ is a nonzero polynomial.  Fixing the first coordinate leaves
boundedly many second coordinates, except for boundedly many full
fibres, each of length $O(T_2)$.  As $T_2\ll T_1$, the count is
$O(T_1)$.  This includes powers of a linear form and needs no uniform
height-free bound for a Thue equation.  Exchanging the pairs handles
bidegree $(a,1)$ with count $O(L_1)$ and right-hand side $hq$.

These cases prove, with arbitrarily small internal losses, the local
bound of \cite[Lemma~8]{reuss2014} in both orientations.  Reducible
auxiliary polynomials have boundedly many factors.  Rational points on
an irreducible factor that is not geometrically integral lie in its
bounded-degree conjugate intersections.  The resulting local count is
therefore uniform in all the variable coefficients.

\medskip\noindent
\textbf{Step 4: the two short-vector summations.}
We give the scale check behind \cite[Section~2.3]{reuss2014}.
For the first lattice, use intervals with
$s_0=(D/E)x_3/M$, $x_3\asymp M$.  A shortest primitive vector
$(x_1,x_2)$ in a box with $L<L_1\leq2L$ satisfies
\[
 x_2x_3=(ME/D)x_1+O(E/L),\qquad
 |x_1|\ll D/L,\quad |x_2|\ll E/L,
 \quad L\gg(DE/M)^{1/2}.
\]
If $x_1x_2\ne0$, the integer product $x_2x_3$ has $O(E/L)$
choices for each of $O(D/L)$ choices of $x_1$.  The divisor bound gives
$N^{o(1)}DE/L^2$ intervals.  Their local contributions sum to
$N^{O(\zeta)}DE/L$; summing the dyadic $L$ gives
$N^{O(\zeta)}(DE M)^{1/2}$.  The primitive axis vectors have one
coordinate equal to $\pm1$ and give the additional costs $D+E$.

For the other lattice choose $U=N/(c_1D^2)$ and $V=N/(c_2E^2)$
within the dyadic comparison constants.  Then
\[
 t_0=(c_1/c_2)(D/E)^2(x_3/M)^2
     =(V/U)(x_3/M)^2.
\]
Replace this by the grid point $(V/U)x_4/M$, where
$x_4=\lfloor x_3^2/M\rfloor$.  The change is within the allowed
interval width.  Since $x_3\asymp M$, this discretization has bounded
multiplicity and $x_4\asymp M$.  In particular no factor $c_1/c_2$
remains.  The same divisor summation gives
$N^{O(\zeta)}\{(UV M)^{1/2}+U+V\}$.
Finally take $\zeta$ and the internal margins sufficiently small,
replace $M$ by $N^\zeta\mathcal M$, and sum the primitive reductions.
This proves \eqref{p1:arith:reussbound}.
\end{proof}

\begin{lemma}[Two small odd kernels]\label{p1:arith:twokernel}
Uniformly for $B\asymp\log N$, fixed $0<\beta<1$, and $\varepsilon>0$,
\[
 G_2(N,N^\beta):=
 \#\{x\in I_N:\#\{n\in V_x:K_B(n)\leq N^\beta\}\geq2\}
 \ll_{\beta,C,\varepsilon}N^{\max(\beta,7/12)+\varepsilon}.
\]
For $7/12\leq\beta<1$ this exponent is sharp in the sense
$G_2(N,N^\beta)=N^{\beta+o(1)}$.
\end{lemma}

\begin{proof}
Two selected integers have the unique forms
$n=c_1d^2u$ and $n+h=c_2e^2v$, where $1\leq h<B$,
$u,v\leq T=N^\beta$, and $c_1,c_2$ are squarefree and $B$-smooth.
There are $N^{o(1)}$ choices for $c_1,c_2,h$.
Fix them and dyadically subdivide $d,e$.  With $D,E$ fixed,
$U\asymp N/(c_1D^2)$ and $V\asymp N/(c_2E^2)$.
The elementary congruence count, for fixed $d,e$, gives in this box
\begin{equation}\label{p1:arith:elementarybox}
 \ll\frac{N|h|}{c_1c_2DE}+DE
 \ll (\log N)^{O(1)}T+DE.
\end{equation}
Indeed the common divisor of the two moduli divides $h$, and
$DE\gg N/(\sqrt{c_1c_2}T)$.
If $c_1c_2(DE)^2<N$, then $DE<\sqrt N$ and
\eqref{p1:arith:elementarybox} suffices.  It also suffices for
$DE\leq N^\delta$, with a small fixed $\delta$.
For $UV\leq N^\delta$ fix $u,v$ and apply
Lemma~\ref{p1:arith:pell}.  These edge regimes cost
$N^{\max(\beta,1/2)+O(\delta)}$.

In the other boxes apply Lemma~\ref{p1:arith:reuss}.  In its notation the
main exponents are
\[
 f_1(z,w)=z/2+\tfrac98z(w-z),\qquad
 f_2(z,w)=w-z+\tfrac98z(w-z).
\]
They meet at $z=2w/3$; $f_1$ is increasing to the left of that point and
$f_2$ decreasing to its right.  Since $0\leq w\leq1+o(1)$,
\[
 \min(f_1,f_2)\leq w/3+w^2/4\leq7/12+o(1).
\]
The remaining terms satisfy $D+E\ll N^{1/2}$ and $U+V\ll T$.
Sum the polylogarithmic boxes and reduce the internal margins.

For the lower bound let $D\asymp\sqrt{N/T}$ with a fixed sufficiently
large comparison constant, and take coprime $d,e\in[D,2D]$.
The conditions $n\equiv0\pmod{d^2}$, $n+1\equiv0\pmod{e^2}$ have
$\gg N/D^4$ solutions in a fixed interior subinterval of $[N,2N)$,
because $\beta>1/2$.  Both remaining cofactors are at most $T$.
The total number of incidences is $\gg D^2N/D^4\asymp T$.
A fixed $n$ is counted at most $\tau(n)\tau(n+1)=N^{o(1)}$ times.
The two neighboring integers lie in a common window, proving the claim.
\end{proof}

\begin{lemma}[A large selected kernel]\label{p1:arith:anchor}
For any fixed $0<\tau<1/2$ and $1\leq T\leq3N$, the number of separated
pairs possessing a mixed square relation with a selected vertex whose
$B$-kernel exceeds $T$ is at most
\[
 N^{3/2+o(1)}T^{-1/2+\tau}.
\]
Also the number of windows containing one vertex with kernel at most $T$
is at most $N^{1/2+o(1)}T^{1/2}$.
\end{lemma}

\begin{proof}
Write $n=a^2cr$, with $c$ squarefree and $B$-smooth and $r=K_B(n)$.
The one-kernel assertion follows from
\[
 \sum_{\substack{n\leq3N\\K_B(n)\leq T}}1
 \leq\sqrt{3N}\prod_{\ell\leq B}(1+\ell^{-1/2})
                       \sum_{r\leq T}r^{-1/2}
 \leq N^{1/2+o(1)}T^{1/2},
\]
with a factor $B$ for its placements.
For the first assertion anchor a selected value $n=x+i$ with $K_B(n)>T$.
Every prime of its kernel must occur at one of the $B$ offsets of the other
window.  There are at most $B^{\omega(r)}$ assignments; each specifies a
residue class modulo $r$.  Since $r\leq3N$, there are $O(N/r)$ possible
other starts, including $r>N$.  Summation over anchors is bounded by
\[
 B N^{3/2}\prod_{\ell\leq B}(1+\ell^{-1/2})
       \sum_{\substack{r>T\ {\rm squarefree}\\\ell\mid r\Rightarrow\ell>B}}
           B^{\omega(r)}r^{-3/2}.
\]
Rankin's inequality bounds the last sum by
\[
 T^{-1/2+\tau}\prod_{\ell>B}(1+B\ell^{-1-\tau})
 =T^{-1/2+\tau}\exp(O_\tau(B^{1-\tau})).
\]
Both products and the placement factor are $N^{o(1)}$.
\end{proof}

\subsection{Multiquadratic surfaces and the curves that can survive}

We use two uniform determinant estimates.  A geometrically integral affine
surface of degree $d$ in $\mathbb A^3$ has, outside
$O_{d,\varepsilon}(H^{2/\sqrt d+\varepsilon})$ integral points of height
at most $H$, a covering by
$O_{d,\varepsilon}(H^{1/\sqrt d+\varepsilon})$ curves of bounded degree
\cite[Lemma~2]{salberger2023}.  An integral curve of degree $e$ in fixed
affine space has $O_{e,n,\varepsilon}(H^{1/e+\varepsilon})$ integral
points in the same box \cite[Theorem~A]{pila1995}.  These constants do
not depend on the defining coefficients.  For a curve irreducible over
$\mathbb Q$ but not geometrically integral, its rational points lie in the
intersection of distinct conjugate components and have bounded cardinality.

Consider equations
\begin{equation}\label{p1:arith:surfaceequations}
 w_\nu^2=d_\nu P_\nu(X)Q_\nu(Y),\quad
 P_\nu(X)=\prod_{i\in I_\nu}(X+i),\quad
 Q_\nu(Y)=\prod_{j\in J_\nu}(Y+j),\quad1\leq\nu\leq k,
\end{equation}
where $I_\nu,J_\nu$ are nonempty, disjoint within each list, and
$|I_\nu|+|J_\nu|\leq K$.  The offsets are $O(B)$ and the nonzero
coefficients $d_\nu$ are squarefree and $O(B)$-smooth.  Constants $k,K$
are fixed.

\begin{lemma}[The uniform geometric counts]\label{p1:arith:geometric}
Fix $A_0,K,\varepsilon>0$.  For sufficiently large fixed $k$, a system
of the form \eqref{p1:arith:surfaceequations}, with arbitrary nonzero integer
coefficients $|d_\nu|\leq N^{A_0}$, has
$O_{A_0,K,\varepsilon}(N^{1/2+\varepsilon})$ integral hosts
$X,Y\asymp N$ at which the even-selection square-class maps on the
two actual sets of selected values are injective and their images are
in direct sum.  If every $Q_\nu$ has degree one, the bound is
$O_{A_0,K,\varepsilon}(N^\varepsilon)$.  If the degrees of the $P_\nu$
vary, select sufficiently many equations of one of their finitely many
possible degrees before applying the assertion.
These estimates are uniform in the coefficients and do not require their
smoothness; any enumeration of such coefficients is a separate cost.
They remain valid after other fixed centers have been specialized into
the $d_\nu$.

For one variable, a bounded-degree system of split-product square equations
with polynomial-height coefficients has $N^{o(1)}$ possible integer values
as soon as its disjoint root sets include at least two distinct selected
offsets.  This includes either one equation of degree at least two or
two equations with disjoint nonempty root sets.

The following graph consequences hold, with the number of component
equations chosen sufficiently large after the stated fixed parameters.
\begin{enumerate}
\item For any $\eta,\varepsilon>0$, pairs of separated nonaligned
logarithmic windows with at least $\eta B$ nontrivial unpinned components
number $O_{\eta,C,\varepsilon}(N^{1/2+\varepsilon})$.
\item For every fixed integer $a\geq1$ and $\varepsilon>0$, there is
$k_a$ such that $k_a$ components having $a$ vertices in one window and
one in the other have $O_{a,C,\varepsilon}(N^\varepsilon)$ dyadic hosts.
For $a=1$ aligned hosts are excluded; for $a\geq2$ no such exclusion
is required.
\item After one start is fixed, a component with at least two vertices
in the other window has $N^{o(1)}$ possible partners, uniformly for a
fixed bound on the component size and polynomial coefficient heights.
\end{enumerate}
\end{lemma}

\paragraph{Parameters in the geometric counts.}
The external estimates are applied in the following embeddings.  Write
$m=\binom{D+2}{2}+k-1$ for the ambient dimension of the affine Veronese chart,
and $q=(a+1+A_0)/2$ in the degree-$(a,1)$ case.
\begin{center}
\small
\begin{tabularx}{\textwidth}{@{}>{\raggedright\arraybackslash}p{.23\textwidth}>{\raggedright\arraybackslash}p{.23\textwidth}>{\raggedright\arraybackslash}p{.15\textwidth}>{\raggedright\arraybackslash}X@{}}
\toprule
External count & Dimension and degree & Point height & Coefficient height\\
\midrule
Surface covering, Steps~1--4
& Surface in $\mathbb A^3$; $2^k\leq d_{\rm im}$,
$d_{\rm im}=O_{k,K}(1)$
& $H\ll N^D$
& Polynomial in $N$; no coefficient-height dependence in the covering constant\\
Lifted curves, Steps~2--4
& Curves in $\mathbb A^m$; $2D\leq e=O_{D,k,K}(1)$
& $H\ll N^D$
& No bound on the coefficients of the covering curves is required\\
Degree-$(a,1)$ surface covering, Steps~5--6
& Surface in $\mathbb A^3$; $2^k\leq d_{\rm im}$,
$d_{\rm im}=O_{k,a}(1)$
& $H\ll N^q$
& Polynomial in $N$; the same coefficient-uniform covering estimate\\
Degree-$(a,1)$ lifted curves, Steps~5--6
& Curves in $\mathbb A^{k+2}$; $E+1\leq e=O_{k,a}(1)$
& $H\ll N^q$
& No coefficient-height bound is required by the curve estimate\\
\bottomrule
\end{tabularx}
\end{center}
The curve rows concern components that can contain the hosts being counted;
Steps~3 and~5--6 exclude the lower degrees.  The input coefficients satisfy
$|d_\nu|\leq N^{A_0}$, and the offsets are $O(\log N)$.
Thus the defining equations and the projected surface equations have
polynomial height with fixed exponents, while the external constants depend
only on the fixed dimensions, degrees and internal error margins.
For the first bound choose $A_0,K,\varepsilon$, then
$D\geq(K+A_0)/2$, then $k$, and finally the internal margin in the
surface and curve counts.  For the degree-$(a,1)$ bound choose $q$, then
$E$ with $q/(E+1)<\varepsilon/8$, then $k$ satisfying the two conditions
in Step~6, and finally that internal margin.  All these choices precede
$N\to\infty$.  The count for one system is uniform over arbitrary
coefficients in the stated height range; the $N^{o(1)}$ enumeration in
the graph application uses smoothness separately.

\begin{proof}
We first count one fixed system, keeping the coefficient height visible;
only afterwards do we sum the patterns arising from the graph.
We will retain a sufficiently large fixed number \(k\) of equations,
chosen after the height parameter \(D\) below.  All degrees and
ambient dimensions below depend on that number and on the fixed parameters,
not on the values of the coefficients.

\medskip\noindent
\textbf{Step 1: geometric integrality and the finite affine projection.}
Over an algebraically closed field of characteristic zero, the square classes on the
right of \eqref{p1:arith:surfaceequations} are independent: the valuation
along a factor $X+i$ detects exactly its own equation.  The coordinate
ring is free of rank $2^k$ over the polynomial ring in $X,Y$, and its
generic localization is a field.  It is therefore integral.

The affine projection $(X,Y,w)\mapsto(X,Y,\sum_\nu w_\nu)$ is finite
and birational onto its image.  Finiteness follows from integrality over
$\mathbb Q[X,Y]$.  Birationality follows because no nontrivial sign change
fixes the sum: the individual radicals are independent basis elements of
the multiquadratic extension.  The image surface in $\mathbb A^3$ has
degree at least $2^k$ and degree bounded in terms of $k,K$.
The projection need not extend across the projective boundary; only its
affine finiteness is used when curves are lifted.

\medskip\noindent
\textbf{Step 2: projective degrees and the height of integral points.}
Choose an integer $D\geq(K+A_0)/2$ before choosing $k$, and homogenize
in the weighted ring
\[
 \mathbb Q[X,Y,Z,W_1,\ldots,W_k]/
 \left(W_\nu^2-d_\nu Z^{2D-|I_\nu|-|J_\nu|}
       \prod_{i\in I_\nu}(X+iZ)\prod_{j\in J_\nu}(Y+jZ)\right),
\]
with $\deg X=\deg Y=\deg Z=1$, $\deg W_\nu=D$.
Its $D$th Veronese ring is generated in degree one by the degree-$D$
monomials in $X,Y,Z$ and the $W_\nu$.  This embeds an integral projective
surface finite of degree $2^k$ over $\mathbb P^2$.  The hyperplane bundle
is the pullback of $\mathcal O_{\mathbb P^2}(D)$, so its degree is
$2^kD^2$ and every integral curve $C$ satisfies
\begin{equation}\label{p1:arith:curvedegree}
 \deg C=D\deg\sigma(C)\,[\mathbb Q(C):\mathbb Q(\sigma(C))].
\end{equation}

Here $\sigma$ is the finite map to $\mathbb P^2$.
At an integral point in the dyadic region,
$|w_\nu|\ll N^{(K+A_0)/2}$.  In the chart $Z=1$, the Veronese
coordinates are integers of height $H\ll N^D$, including the $W_\nu$.
The same bound holds for the affine projection in Step~1.  Thus both
applications of the determinant estimates use one common height bound.

\medskip\noindent
\textbf{Step 3: excluding the curves of smallest possible degree.}
A curve of degree $D$ maps with degree one to a line.  A vertical or
horizontal line meeting the positive region is impossible: a nonconstant
product of distinct linear factors in the other variable would have to be
a rational square.  On an oblique line, square products force every root
on one side to match exactly one root on the other.  Two component
equations give two distinct exact units, hence a positive rational
alignment of height $O(B)$.  Thus a nonaligned point is on no curve of
degree $D$; by \eqref{p1:arith:curvedegree}, its curves have degree at least
$2D$.

In the direct-sum formulation the
exceptional degree-$D$ curves are excluded for a more intrinsic reason:
their two exact units give nonempty even selections, two values on each
side, with equal square classes.  Injectivity makes both classes nonzero,
contradicting the direct sum.  Thus the argument does not require absence
of geometric alignment when the actual spaces are already in direct sum.

\medskip\noindent
\textbf{Step 4: the surface covering and the first graph consequence.}
Apply the surface covering, lift its curves by the finite affine map, and
take their closures in this embedding.  Their numbers of components and
degrees are bounded in terms of the fixed parameters by B\'ezout.
Pila's estimate and the degree lower bound $2D$ give, per pattern,
\[
 O\big(N^{2D2^{-k/2}+\varepsilon'}
        +N^{1/2+D2^{-k/2}+\varepsilon'}\big).
\]
Choose $k$ large enough after $D,\varepsilon$ and reduce the internal
margin $\varepsilon'$.  This proves the per-system bound.

For the graph consequence, smoothness gives $|d_\nu|\leq N^{A_0}$ for a
fixed $A_0$.  At least half the components in a population $c\geq\eta B$ have size
at most $K=\lceil4/\eta\rceil$.  Select $k$ of them.
The offsets and smooth coefficients have
$B^{O_K(k)}2^{k\pi(O(B))}=N^{o(1)}$ patterns.  Choosing $k$ large after
$D,\varepsilon$, and summing these patterns, proves the first assertion.

\medskip\noindent
\textbf{Step 5: reduced fibres for equations of type $(a,1)$.}
For the stronger degree-one assertion we exclude all curves of a fixed
degree at most $E$, before making $k$ large.  For this part use the
ordinary affine embedding with coordinates $X,Y,w_1,\ldots,w_k$.
Write the equations as
$w_\nu^2=d_\nu P_\nu(X)(Y+j_\nu)$ with $\deg P_\nu=a$.
On the normalization of an integral curve of degree $e\leq E$,
the coordinate functions $X,Y$ have degrees $m_X,m_Y\leq e$ and the
genus is at most $(e-1)(e-2)/2$.  Neither function can be constant:
otherwise the many distinct roots in the other variable would all be
branch values, contradicting Riemann--Hurwitz.
Discard equations having a branch value among their selected roots, or a
finite fibre meeting the pole divisor of the other coordinate.
At most
\[
 (2g-2+2m_X)+(2g-2+2m_Y)+m_X+m_Y\leq2e^2
\]
equations are lost.  For each remaining equation all finite fibres are
reduced and avoid the poles.  Parity of divisors gives
\[
 F_Y(-j_\nu)=\sum_{i\in I_\nu}F_X(-i),\qquad m_Y=am_X.
\]
In fact
\[
 \operatorname{div}\!\left(\frac{P_\nu(X)}{Y+j_\nu}\right)
       =F_Y(\infty)-aF_X(\infty),
\]
independently of $\nu$.  For two good equations the quotients therefore
differ by a nonzero constant.  Eliminating this common quotient gives
$Y=R(X)$ with $\deg R\leq a$.  Since
$m_Y=(\deg R)m_X=am_X$, the degree is exactly $a$.

\medskip\noindent
\textbf{Step 6: interpolation on the offsets and exclusion of low-degree curves.}  If
$k\geq4(E+1)^2$, there are at least three good equations.
Their $3a$ distinct grid points $(-i,-j_\nu)$ lie on the graph of $R$:
the reduced-fibre identity shows that $R$ is finite and equals $-j_\nu$
at each selected $-i$.  These grid points determine $R$.
Choose $2a+1$ of them.  Rational interpolation, using maximal minors of
the integer evaluation matrix, gives coprime integer polynomials
$P,Q$ of degrees at most $a$, coefficients $B^{O_a(1)}$, and $R=P/Q$.
The kernel of this interpolation system is one-dimensional: two solutions
would have cross-product of degree at most $2a$ vanishing at $2a+1$
distinct abscissae.  Since the reduced rational function has degree exactly
$a$, no nonconstant polynomial multiple of its coprime numerator and
denominator still has both degrees at most $a$.
This makes the coefficient uniformity explicit: the interpolation data
are the offsets, so their height bound is independent of the $d_\nu$.

If $Q$ is nonconstant, the nonzero integer resultant
$\operatorname{Res}(P,Q)$ has size $B^{O_a(1)}$ and $Q(X)$ divides it
at an integral point.  Thus $|X|\leq B^{O_a(1)}$.  If $Q$ is constant
and $a\geq2$, then $|Y|\gg |X|^a/B^{O_a(1)}$, impossible for
$X,Y\asymp N$.  If $a=1$ and $Q$ is constant, the graph is a rational
line; its grid points give an exact positive alignment.  Thus every
remaining dyadic nonaligned point lies on a curve of degree greater than
$E$.

In the ordinary affine embedding the height is $H\ll N^q$ for a fixed
$q=(a+1+A_0)/2$.  Choose $E$ so that $q/(E+1)<\varepsilon/8$, then
$k\geq4(E+1)^2$ so that $2q2^{-k/2}<\varepsilon/8$.
The same finite projection, surface covering and curve estimate now give
$N^\varepsilon$ per fixed coefficient pattern.  For $a=1$, its exceptional
line is also excluded under the direct-sum hypothesis: two grid matches
give the forbidden even selections on the two actual blocks.  Hence the
degree-one per-system estimate is proved with arbitrary polynomial-height
coefficients.  In the graph application one sums the $N^{o(1)}$ smooth
patterns, proving the second graph assertion, including $a=1$.

\medskip\noindent
\textbf{Step 7: fixing one center and reducing to a quadratic equation.}
For the last assertion fix the first start and the component data.  The
other start satisfies $\prod_{j\in J}(Y+j)=ev^2$, with fixed squarefree
$e\leq N^{O(1)}$ and $|J|\geq2$.  Outside primes dividing
$e\prod_{i\ne j}(i-j)$, every selected factor has even valuation.
Fixing the squarefree parts of two factors costs
$2^{O(\omega(e\prod_{i\ne j}(i-j)))}=N^{o(1)}$ for fixed $|J|$.
Their difference is a nonzero fixed offset, so
Lemma~\ref{p1:arith:pell} gives $N^{o(1)}$ starts.  This is a
bounded-height argument; no height-free theorem for all integral points
is being invoked.
The one-variable assertion with two separate equations follows in exactly
the same localization: take one selected factor from each equation and
put the prime divisors of both coefficients and both within-equation
discriminants into the exceptional prime set.  Outside it each chosen
factor has even valuation.  The two distinct offsets again give one
quadratic equation counted by Lemma~\ref{p1:arith:pell}.
\end{proof}

\begin{lemma}[The dense weighted sector]\label{p1:arith:dense}
For every fixed $\eta,\varepsilon>0$,
\[
 \sum_{\substack{x,y\in I_N,\ |x-y|>L\\
        (x,y)\ {\rm nonaligned},\ c(x,y)\geq\eta B}}
       (2^{\rho(x,y)}-1)
 \ll_{C,\eta,\varepsilon}N^{1+\varepsilon}.
\]
The same bound holds for the full value kernel, without block-parity rows.
\end{lemma}

\begin{proof}
Let $t_2,t_3$ count nontrivial unpinned components with two and three
vertices.  If either exceeds a sufficiently large fixed threshold,
Lemma~\ref{p1:arith:geometric}(2), in both orientations for size three,
gives $N^\varepsilon$ hosts.  Their weight is at most
$2^{B+o(B)}=N^{1+o(1)}$, since each other component needs two vertices
and the number of defective vertices is $o(B)$.
Otherwise $t_2,t_3=O(1)$ and the vertex budget gives
$4c-2t_2-t_3\leq2B$, hence $c\leq B/2+O(1)$.
Lemma~\ref{p1:arith:geometric}(1) gives $N^{1/2+\varepsilon}$ hosts,
each of weight at most $2^{B/2+o(B)}=N^{1/2+o(1)}$.
Reduce the internal margins in both estimates.
\end{proof}

\begin{lemma}[Uniform residual rank on an aligned pair]\label{p1:arith:alignedrank}
For separated aligned windows with reduced height
$b=\max(a,b)\leq B^3$, let $u$ be the number of exact units.  Then
\[
 \rho(x,y)\leq(u-1)_++o(B)\leq L/b+o(B),
\]
uniformly over the channels. For fixed channel height $b$ and fixed
$D\ge6b$, an enlarged prime cutoff $6bB\le H\le DB$ leaves
residual $O_{C,b,D}(B/\log B)$.
\end{lemma}

\begin{proof}
Normalize the two windows to $bV_x$ and $aV_y$, and merge equal values.
All remaining distinct integers lie in an interval of diameter $O(B^4)$
centered at an integer at least $N$.  Remove prime rows $\ell\leq B$
and rows with $\ell\mid ab$.  Every other prime row meets at most two
distinct normalized vertices, since it occurred at most once in each
original window.  Its equations again leave one variable per unpinned
component.  Let $c_K$ count these components of size at most a fixed $K$.

For these components, the removed prime rows and one total-parity row
form a binary matrix of redundancy at most
$s=\pi(B)+\omega(ab)+1=O(B/\log B)$.  A nonzero kernel word gives
an even square product of distinct normalized integers.  By
Lemma~\ref{p1:arith:runge}, it uses at least
$\delta\gg\log N/\log B$ vertices, and hence at least $\delta/K$
components.  The Hamming argument in Lemma~\ref{p1:arith:defects} yields
$c_K=O_{C,K}(B/\log B)$.  Components larger than $K$ number at most
$2B/K$.  Since $K$ can be chosen arbitrarily large before $N$, the
dimension of the value kernel on the distinct normalized set is $o(B)$.

The kernel of the merging map on the original relative selections is
exactly the even-selection space of the $u$ exact units, of dimension
$(u-1)_+$.  Its image lies in the value kernel just bounded, proving
the first assertion.  For fixed height the normalized interval has
diameter $O(B)$.  Private primes above a sufficiently large fixed multiple
of $B$ leave only $O(B/\log B)$ defective normalized vertices, directly
by Lemma~\ref{p1:arith:defects}; this gives the last assertion.
\end{proof}

\subsection{Asymmetric extensions and homogeneous seeds}

\begin{lemma}[An asymmetric pair profile]\label{p1:arith:asymmetric}
Fix integers $D,Q\geq1$ and $D_{\rm cut}\geq2D+1$.
Let $B=L+1$ and $|L-\log N/\log2|\leq C$.
For each $N\leq X<N+QN$, let $U_X$ be any subset of
$X+[-DB,DB]$, and let $V_z$ be a new window with $z\in I_N$.
Choose an integer cutoff
$(2D+1)B\leq H\leq D_{\rm cut}B$.
For this asymmetric comparison call $(X,z)$ aligned at radius $D$ if
there are coprime positive integers $r,s\leq B^3$ such that
$|rz-sX|\leq D(r+s)B$; use its complement in the sum below.
Let $\delta^*(X,z)$ be the dimension of the even selections
of $V_z$ whose prime valuations above $H$ belong to the value space of
$U_X$.  Then
\[
 \sum_{\substack{N\leq X<N+QN,\ z\in I_N\\
                 (X,z)\ {\rm nonaligned\ at\ radius}\ D}}
       (2^{\delta^*(X,z)}-1)
 \ll_{C,D,D_{\rm cut},Q,\varepsilon}N^{31/24+\varepsilon}.
\]
The choice $Q=1$ gives the two-dyadic-window version.
Adding prime coordinates at most $H$ to the old data leaves
$\delta^*$ unchanged.
\end{lemma}

\begin{proof}
All alignments in this proof use the radius $DB$ of the old enclosing
interval.  The cutoff lower bound is larger than both interval diameters,
so a prime above $H$ occurs at most once in each block.
The fixed expanded old-centre range changes only constants in the
biform support estimates and the bounded-degree graph counts.
For the graph above $H$, let $D_z$ be the number of new defective vertices
and let $c$ count nontrivial unpinned components.  Projection to the new
block gives $\delta^*\leq\min(L,D_z+c)$, with $D_z=o(B)$.
A nonzero even selection uses at least two new vertices.  If it uses no
old vertex, the new window has two defects and contributes $N^{1+o(1)}$
hosts after the old parameter is chosen.  Otherwise it uses at least one
old and two new vertices.  Set $T=N^{5/12}$.
If all their kernels are at most $T$, Lemmas~\ref{p1:arith:twokernel}
and~\ref{p1:arith:anchor} give
\[
 N^{7/12+\varepsilon'}N^{1/2+5/24+o(1)}=N^{31/24+\varepsilon'+o(1)}.
\]
If one kernel exceeds $T$, the same CRT anchor proof gives
$N^{3/2+o(1)}T^{-1/2+\tau}=N^{31/24+(5/12)\tau+o(1)}$.
All these arguments concern the large-prime valuations only, so they do
not impose parity on the old selection.  They are unchanged at the
fixed enlarged threshold $H$.

For $c<\eta B$, the weight is at most $2^{\eta B+o(B)}$, and the support
bound suffices.  For $c\geq\eta B$, the proof of
Lemma~\ref{p1:arith:geometric}(1) applies to the unequal fixed widths and
gives $N^{1/2+\varepsilon'}$ hosts.  Fix a large integer $K$.  Components
with one new vertex and more than $K$ old vertices number $O_D(B/K)$.
A sufficiently large fixed population of any other type $(a,1)$,
$1\leq a\leq K$, has only $N^{\varepsilon'}$ hosts by
Lemma~\ref{p1:arith:geometric}(2).  On this exceptional set use
$2^{\delta^*}\leq2^L\ll N$.  Else the number $c_1$ of one-new-vertex
components is $O_D(B/K)+O_K(1)$, and $2c-c_1\leq B$.  Thus
$\delta^*\leq B/2+O_D(B/K)+o(B)$ and the remaining dense mass is
$N^{1+O_D(1/K)+\varepsilon'}$.
Choose $K$, then $\eta,\tau,\varepsilon'$ after the required margin.
Finally shrinking the old set only decreases $\delta^*$, and prime bits
at most $H$ have already been forgotten.
\end{proof}

\paragraph{Alignment convention for asymmetric applications.}
The enlarged radius in this lemma follows the actual old interval.
Original pairs of equal windows retain the alignment convention in
Section~\ref{p1:arith:section}.  Whenever a source is compressed or an enclosing
interval is enlarged below, its aligned sector is enlarged at the
same time; that sector still has $NB^{O(1)}$ hosts and is treated by
the separate aligned-rank estimates.  The radius-one assertion with
arbitrary $D$ is false, as the following family shows.

\paragraph{Necessity of the enlarged alignment radius.}
In Lemma~\ref{p1:arith:asymmetric}, the alignment radius follows the
old enclosing interval $X+[-DB,DB]$: the exclusion is
$|rz-sX|\le D(r+s)B$. The centres satisfy
$N\le X<N+QN$, $z\in I_N$, with fixed $Q\ge1$, and the cutoff satisfies
$(2D+1)B\le H\le D_{\rm cut}B$. The weighted bound for $2^{\delta^*}-1$ has power $31/24$.
A shifted $2:3$ family with
$L=3k$, $B=3k+1$, $X=2t$, $z=3t+3B$ and
$U_X=[X,X+3B]\cap\Z$ has $\delta^*\ge k-1$. It is nonaligned at
radius one for $t>5B^4$ but aligned at radius two; its dyadic mass
of order $N^{4/3}$ refutes a radius-one estimate for arbitrary $D$.
Compressed-source applications use the same enlarged radius on both
sides of the alignment partition.

\begin{lemma}[Aligned loss above a rational pair]\label{p1:arith:palmrank}
Fix a geometrically good pair in channel $a:b=2:3$ or $3:4$, regardless
of its affine compatibility.  Among targets $z$ with
$|z-x|,|z-y|\geq L$ and distinct from the two sources, there are
$B^{O(1)}$ aligned targets per source.  Except for $B^{O(1)}$ source
pairs at dyadic endpoints, their additional relative rank loss modulo
the old value space is at most
\[
 L/3+o(L)\quad(a:b=2:3),\qquad
 L/2+o(L)\quad(a:b=3:4).
\]
The result allows the prime bits of $a,b$ in the old information.
\end{lemma}

\begin{proof}
Represent the source by $x=at-i_0$, $y=bt-j_0$ and put $X=abt$.
The true old normalized values form a subset of
$X+[-bB,bB]$.  Applying the residual argument of
Lemma~\ref{p1:arith:alignedrank} to a fixed-width enclosing interval and
the new interval leaves $o(B)$ dimensions after exact units are removed.
If an exact unit of the enlarged graph uses an old value absent from the
true set, its new nondefective vertex is pinned in the true graph.
Consequently the remaining loss is at most the number of new vertices
matched to actual old values, plus $o(B)$.

Write an aligned target as $z=(ct+h)/d$, with $c,d$ positive coprime,
$c,d\leq B^{O(1)}$ and $h/d=O(B)$.  A match to the first old family
requires
$h+dk\equiv0\pmod{c/(c,a)}$; the second family requires the same
condition modulo $c/(c,b)$.  Their intersection is the class modulo $c$,
because $(a,b)=1$.  Thus the union has at most
\begin{equation}\label{p1:arith:matchdensity}
 B\frac{(c,a)+(c,b)-1}{c}+O(1)
\end{equation}
new offsets.  Conditions on the true intervals can only reduce this count.

For $2:3$, the source range is
$N/2-O(B)\leq t\leq2N/3+O(B)$, so $3/2\leq c/d\leq4$ apart
from endpoint errors of order $B/N$.  A distinct rational of polynomial
height cannot cross either endpoint by such a small amount.
For $c\geq7$ the density in \eqref{p1:arith:matchdensity} is at most
$1/3$ (check $7\leq c\leq12$, then bound the numerator by four).
For $c=5$ it is $1/5$.  The cases $c/d=2,3$ physically reproduce one
old family.  Their supports intersect the new support in at most one
integer, because the starts are at distance at least $L$.  The other
family supplies at most $B/3+O(1)$ units, counted in its actual window
of width $B$.  The only remaining possibilities with $c\leq6$ are
$c/d=4,3/2$; these force $t=N/2+O(B)$ or $t=2N/3+O(B)$.
They account for only $B^{O(1)}$ source pairs.

For $3:4$, one has $N/3-O(B)\leq t\leq N/2+O(B)$ and
$2\leq c/d\leq6$.  The density is at most $1/2$ for $c\geq9$
(check $9\leq c\leq12$, then use numerator at most six), and also
for $c=5,7,8$.  The physical cases $c/d=3,4$ leave at most one unit
in one old family and $B/4+O(1)$ in the other.  The remaining fractions
$2,6$ force $t=N/2+O(B)$ or $N/3+O(B)$ and are endpoint sources.
Finally the canonical slopes and displacements have polynomially many
choices in $B$, proving the target count.
\end{proof}

\begin{proposition}[Homogeneous seeds for fixed extensions]
\label{p1:arith:fixed-homogeneous-seeds}
Let $\mathcal A_{\rm raw}$ be the geometric good-edge family obtained by
omitting the compatibility condition in the definition of $\mathcal E_N$.
For a finite set $I\subset G_N$ put
\[
 W(I)=p^{|I|}2^{\rho(I)}.
\]
All configurations in the following sums have distinct, pairwise separated
starts.  For every $\varepsilon>0$:
\begin{enumerate}
\item The sum of $W(\{x,y\})$ over dependent pairs outside
$\mathcal A_{\rm raw}$ is $O_{C,\varepsilon}(N^{-4/5+\varepsilon})$.
Moreover $\sum_{e\in\mathcal A_{\rm raw}}W(e)\ll_C LN^{-2/3}$.
\item For a geometric good source $e\in\mathcal A_{\rm raw}$ in channel
$a:b\in\{2:3,3:4\}$, sum over targets $z$ for which
$S_z\cap\sum_{x\in e}S_x\ne\{0\}$. Then
\[
 \sum_{e,z}W(e\cup\{z\})\ll_{C,\varepsilon}N^{-41/30+\varepsilon}.
\]
\item For two distinct geometric good edges, whose union has three or
four starts,
\[
 \sum_{e\ne f}W(e\cup f)
 \ll_C L^2N^{-4/3}.
\]
\end{enumerate}
These are rank bounds without any affine compatibility restriction.
\end{proposition}

\begin{proof}
The homogeneous indicators
$H_x=\ind_{\{f(x-1)=\cdots=f(x+L-1)\}}$ have zero affine target, so
$\E\prod_{x\in I}H_x=W(I)$. The first assertion follows from the
positive profile and the exceptional-channel estimates in
Theorem~\ref{p1:arith:pairs}; on a dependent pair,
$2^\rho\le2(2^\rho-1)$.

For (2), the good source is dependent, so positive crossing gives a
connected triple.
The separated homogeneous variant in
Section~\ref{p14:comp:third:aligned-family}, with at most three marked
source pairs per triple, gives $O(N^{-41/30+\varepsilon})$.

\paragraph{Two interacting raw edges.}
Write $W_e=\sum_{x\in e}S_x$ and sum over ordered pairs:
\[
 \mathfrak I_N=
 \sum_{\substack{e,f\in\mathcal A_{\rm raw},\ e\ne f\\
                  e\cup f\ {\rm separated},\ W_e\cap W_f\ne\{0\}}}W(e\cup f).
\]
We prove $\mathfrak I_N\ll_{C,\varepsilon}N^{-41/30+\varepsilon}$.
If the two edges share a vertex, their union is a connected triple;
the same homogeneous estimate, with bounded marking multiplicity, applies.

For disjoint edges first treat nonaligned compressed blocks, using
a fixed radius enclosing both true value sets in the alignment test.
Its complementary enlarged aligned sector is retained below.  Write
$r_f=|U_f|-1\geq5L/3-O(1)$ and let $\delta^*$ be the loss of the
even-selection rank on $U_f$ after the old normalized values and small
primes have been exposed.  The excess above $W(e)W(f)$ is at most
$W(e)W(f)(2^{\delta^*}-1)$.
The weak-component support proof of Lemma~\ref{p1:arith:asymmetric}
still applies to two subsets of fixed-width logarithmic intervals:
it needs two selected new values and one selected old value.
It gives $N^{31/24+\varepsilon}$ hosts.  In this sector
$W(e)W(f)\leq N^{-10/3+o(1)}$, so its mass is
$N^{-49/24+\varepsilon}$.

In the dense sector, the multiquadratic argument gives
$N^{1/2+\varepsilon}$ hosts.  A fixed large population of components
with one new vertex and boundedly many old vertices has only
$N^\varepsilon$ hosts, by Lemma~\ref{p1:arith:geometric}(2).
On that exceptional population use $W(f)2^{\delta^*}\leq1$ and
$W(e)\leq N^{-5/3+o(1)}$.  Otherwise all but
$O(B/K)+O_K(1)$ components use two new vertices.  Hence
\[
 \delta^*\leq|U_f|/2+O(B/K)+o(B),\qquad
 W(f)2^{\delta^*}\leq N^{-5/6+O(1/K)+o(1)}.
\]
The remaining dense mass is $N^{-2+\varepsilon}$.
The entire nonaligned excess is therefore $N^{-5/3+\varepsilon}$.
On an interacting pair let $\delta=\dim(W_e\cap W_f)>0$.
Since $2^\delta\le2(2^\delta-1)\le2(2^{\delta^*}-1)$, the same
nonaligned excess bound controls its full homogeneous mass.

For disjoint edges with aligned compressed centres, the fixed channel,
offset and alignment codes leave $NB^{O(1)}$ quadruples. Choose $z\in f$.
If $S_z\cap W_e\ne0$, the triple $e\cup\{z\}$ is connected.
For each such triple the mate of $z$ in $f$ has $O(B)$ choices,
because the raw good-edge graph has degree $O(B)$. Since
$W(e\cup f)\le W(e\cup\{z\})$, the homogeneous triple bound applies.
If $S_z\cap W_e=0$, instead
$W(e\cup f)\le W(e\cup\{z\})=W(e)p\le N^{-\kappa_b-1+o(1)}$.
The host count gives $N^{-\kappa_b+o(1)}$ in this case.
This includes a positive intersection between the two pair spaces
even when both individual target crossings vanish.
Thus all sectors satisfy the asserted bound for $\mathfrak I_N$.

Finally, pairs with disjoint effective spaces factor exactly. Hence
\[
 \sum_{\substack{e,f\in\mathcal A_{\rm raw},\ e\ne f\\e\cup f\ {\rm separated}}}
       W(e\cup f)
 \le \Bigl(\sum_{e\in\mathcal A_{\rm raw}}W(e)\Bigr)^2+\mathfrak I_N.
\]
The first assertion bounds the squared sum by $O_C(L^2N^{-4/3})$.
Fix an internal margin below $1/30$; since $41/30>4/3$,
the interaction term is absorbed into that scale with a fixed
constant. This proves (3).
\end{proof}

\subsection{Support and branching at a fixed number of blocks}

\begin{lemma}[Small-kernel counts for fixed-width blocks]
\label{p1:arith:small-kernels}
Let $U_X\subset X+[-DB,DB]$ with $D$ fixed and $X\asymp N$, and let
$H\asymp B$ exceed the diameter.  Uniformly for $1\leq T\leq CN$,
the number of translations containing one value with $K_H(n)\leq T$
is at most $N^{1/2+o(1)}T^{1/2}$.  For fixed $0<\theta<1$, the number
containing two such distinct values with $T=N^\theta$ is
$O_{C,D,\theta,\varepsilon}(N^{\max(\theta,7/12)+\varepsilon})$.
In particular it is $O(N^{\theta+\varepsilon})$ for $\theta\geq7/12$.
\end{lemma}

\begin{proof}
Enlarge the selected set to its interval.  The proofs of
Lemmas~\ref{p1:arith:anchor} and~\ref{p1:arith:twokernel} use only a
fixed logarithmic width and a prime threshold of that order.  In particular
$2^{\pi(H)}=N^{o(1)}$, the smooth parts have fixed polynomial height,
and all changes of offset count are powers of $B$.
\end{proof}

Here a projection is an element of the effective even-selection image
in the square-class space, rather than a formal coordinate vector.

\begin{proposition}[A relative relation meeting several blocks]
\label{p1:arith:fixed-support}
Fix integers $r\geq s\geq2$.  For $r$ separated logarithmic windows,
the number of tuples having a relative square relation whose projection
is nonzero in at least $s$ blocks is
\[
 O_{C,r,\varepsilon}\left(N^{r-s/2+\varepsilon}\right).
\]
The same assertion holds for effective even-selection images on subsets
of fixed-width logarithmic intervals, including the compressed good edges,
with only $B^{O_r(1)}$ choices of fixed offsets and fixed dilation factors.
\end{proposition}

\begin{proof}
Choose $s$ blocks on which a witnessing relation has nonzero effective
projection, and fix the other $r-s$ centres. Their values generate the
old space in Theorem~\ref{p13:half:conditional}. The selected $s$
nonzero projections sum into that old space; they need not remain
nonzero modulo it. The theorem therefore gives
$O(N^{s/2+\varepsilon'})$ choices for the selected centres.
Summing the $O_r(1)$ choices of blocks and $O(N^{r-s})$ other centres
gives $O(N^{r-s/2+\varepsilon'})$.

The same theorem permits the effective images of all fixed-width blocks
in its stated height range. Fixed dilations preserve that range, and
the $B^{O_r(1)}$ offset and motif choices are absorbed by reducing
$\varepsilon'$. Injectivity of the local selection maps is unnecessary:
the projections must be nonzero in their own effective images, even when
they vanish modulo the old space. Formal vectors with zero local image
are not counted.
\end{proof}

\begin{proposition}[A fixed number of exposed source windows]
\label{p1:arith:fixed-branching}
Fix $j\geq1$ and source starts $x_1,\ldots,x_j\in I_N$.
Outside a deterministic set of $N^{1/2+o_{C,j}(1)}$ possible starts,
the new relative vector is uniform and independent of every value in
the old windows.  An explicit bound for the exceptional set is
\[
 j(2L+1)+B\sqrt{3N}\exp(2\sqrt{m_{j,N}}),\qquad
 m_{j,N}=\left\lceil\pi(B)+\frac{jB\log(3N)}{\log B}\right\rceil.
\]
Fixed dilations, subsets, and fixed enlarged logarithmic windows satisfy
the same $N^{1/2+o(1)}$ conclusion.
\end{proposition}

\begin{proof}
Let $\mathcal P$ contain all primes at most $B$ and all primes dividing
an old value.  It has at most $m_{j,N}$ members.
The number of integers at most $3N$ whose squarefree part is supported
on $\mathcal P$ is at most
\[
 \sqrt{3N}\prod_{\ell\in\mathcal P}(1+\ell^{-1/2})
 \leq\sqrt{3N}\exp(2\sqrt{|\mathcal P|}),
\]
because the $i$th distinct prime is at least $i$.
Discard overlapping starts and windows containing two such values;
the stated estimate follows even by charging each such window to one
of its values and multiplying by $B$.
If a nonzero new relative combination belonged to the old value space,
its even selection would contain at least two integers.  Every large
odd prime of each selected integer is private inside the new window,
so it would have to divide an old value.  Both selected integers would
have squarefree part supported on $\mathcal P$, contrary to the retained
set.  The new rank therefore increases by exactly $L$ over the old
value space, which proves the asserted conditional uniformity.
For fixed $j$, $m_{j,N}=O_{C,j}((\log N)^2/\log\log N)$.
\end{proof}

\subsection{Products grouped by their ternary support}

\begin{proposition}[A ternary group budget]\label{p1:arith:ternary-budget}
Let $U_i\subset X_i+[-DB,DB]$, $X_i\asymp N$, be three sets with
injective even-selection square-class maps whose images are pairwise
in direct sum.  The width constant and the number of offset patterns
are fixed.  Choose $H$ greater than the three diameters.  For an offset
triple $u=(u_1,u_2,u_3)$ let
\[
 G_u=\prod_{\substack{\ell>H\ {\rm prime}\\
       v_\ell(X_i+u_i)\ {\rm odd}\ (i=1,2,3)}}\ell.
\]
Retain the groups greater than one and write their decreasing values
$G_{(1)}>G_{(2)}>\cdots$.  Put
$\mathcal G_k=\prod_{i=1}^kG_{(i)}$, or zero if fewer than $k$ exist.
For every fixed $k\geq3$, $\varepsilon>0$, and $1\leq T\leq N$,
\begin{equation}\label{p1:arith:ternarycount}
 \#\{\text{tuples}:\mathcal G_k>T\}
 \ll_{C,D,k,\varepsilon}N^\varepsilon
      (N^3/T^2+N^{6/5}).
\end{equation}
\end{proposition}

\begin{proof}
A prime larger than $H$ has only one occurrence in each block, so the
groups are pairwise coprime and
$G_u\mid\gcd(X_1+u_1,X_2+u_2,X_3+u_3)$.
Two groups satisfy $G_uG_v\ll NB$.
If they share an occurrence this follows from divisibility of that
integer.  Otherwise choose two blocks in which the offsets differ.
Their product divides a determinant $nm'-n'm$ of size $O(NB)$.
Its vanishing would give a nonzero relative relation on those two
blocks, contrary to their direct sum.

Three groups with product greater than a sufficiently large constant
times $NB^2$ force their three value vectors to have zero determinant.
The corresponding offset triples cannot be collinear.  Indeed, writing
$u_j=u_1+t_jd$, with $d$ a primitive integer vector and $t_j$ integers,
in the collinear case, all three coprime groups divide
$d\times(X+u_1)$.  If this vector is nonzero, its size $O(NB)$ gives a
contradiction.  If it is zero, two positive value vectors are proportional
with distinct coordinates in every block, again giving a forbidden pair
relation.  Thus the centers lie on a genuine plane
\[
 ((u_2-u_1)\times(u_3-u_1))\cdot X
 =-((u_2-u_1)\times(u_3-u_1))\cdot u_1.
\]
There are only $B^{O(1)}$ such planes, determined by the offsets alone.

For fixed $k$ group first by the squarefree product $Q=\mathcal G_k$.
The ordered factorizations cost $d_k(Q)\ll_{k,\varepsilon'}Q^{\varepsilon'}$,
and the supports cost $B^{3k}$.  CRT fixes one residue class modulo $Q$
for each center.  Summing $O((N/Q+1)^3)$ gives, for
$1\leq T\leq N\leq R\leq N^{O(1)}$,
\begin{equation}\label{p1:arith:groupcrt}
 \#\{T<\mathcal G_k\leq R\}
 \ll_{k,\varepsilon}N^\varepsilon(N^3/T^2+R).
\end{equation}
For $\mathcal G_3>R\gg NB^2$, use its forced plane.  The two largest
groups have product $P>R^{2/3}$ and $P\ll NB$.  Two free coordinates
are fixed by CRT modulo $P$; the plane determines the third.
Summing over the two positive group sizes and their offsets gives
\begin{equation}\label{p1:arith:planecrt}
 \#\{\mathcal G_3>R\}
 \ll_\varepsilon N^\varepsilon(N^2R^{-2/3}+N).
\end{equation}
Here one uses
$\sum_{ab>A}(ab)^{-2}\ll\log(2A)/A$,
$\sum_{ab\leq U}(ab)^{-1}\ll\log^2(2U)$, and
$\#\{ab\leq U\}\ll U\log(2U)$.
The third numerical group is not enumerated again: it certifies the
plane and the lower bound for $P$.

For $k=3$, take $R=N^{6/5}$ in
\eqref{p1:arith:groupcrt}--\eqref{p1:arith:planecrt}.
For induction on fixed $k$, the region
$T<\mathcal G_k\leq N^{6/5}$ is already covered by
\eqref{p1:arith:groupcrt}.  Beyond it,
$\mathcal G_{k-1}>N^{6(k-1)/(5k)}$.
At $k=4$ the preceding $N^3/T^2$ term is $N^{6/5}$; at $k=5$
it is smaller; for $k\geq6$ the threshold is at least $N$ and the
preceding tail at $N$ suffices.  Only a fixed number of induction steps
is taken.  This proves \eqref{p1:arith:ternarycount}.
\end{proof}

\subsection{The probability of two active good edges}

\begin{lemma}[A positive second moment for good edges]
\label{p1:arith:twogoodpositive}
Let $\Xi_N=\sum_{e\in\mathcal E_N}\prod_{x\in e}J_x$.
Then
\[
 \E(\Xi_N)_2=\beta_N^2+o_C(\beta_N^2)
 \ll_C L^2N^{-4/3}.
\]
In particular
$0\leq\beta_N-\P(\Xi_N\geq1)\leq\tfrac12\E(\Xi_N)_2$.
\end{lemma}

\begin{proof}
For compatible good edges put $q_e=q(e)=W(e)$ and
$W_e=\sum_{x\in e}S_x$. On pairs with $W_e\cap W_f=\{0\}$,
their original affine events factor exactly.
On an interacting separated pair,
\[
 |q(e\cup f)-q_eq_f|\le W(e\cup f):
\]
the joint is zero or $W(e\cup f)$, while $q_eq_f\le W(e\cup f)$.
The interaction estimate proved in
Proposition~\ref{p1:arith:fixed-homogeneous-seeds} therefore bounds
the total separated covariance by $O(N^{-41/30+\varepsilon})$.

The diagonal and strict-contact baseline products have only
$NB^{O(1)}$ hosts, using the geometric degree $O(B)$.
Each product is at most $N^{-10/3+o(1)}$, giving a total
$O(N^{-7/3+o(1)})$;
strict contacts have zero original joint probability. Thus
\[
 \E(\Xi_N)_2=\beta_N^2+O_{C,\varepsilon}(N^{-41/30+\varepsilon}).
\]
Choose $\varepsilon<1/30$ and use $\beta_N\asymp_C LN^{-2/3}$.
This gives the relative asymptotic and its fixed-constant upper bound.
Finally $(m-1)_+\le m(m-1)/2$ for integers $m\ge0$ gives the
stated Bonferroni inequality after expectation.
\end{proof}


\subsection{Local rank at critical length}

We apply the binary Griesmer bound
\cite[Section~2, equation~(2.10)]{griesmer1960} to the code of
local even square relations. The arithmetic input is its minimum
nonzero weight; the residual-code proof below keeps the dependence
on the parameters explicit.

\begin{lemma}[Local rigidity]
\label{p1:limits:local-rigidity}
Let $V=\{U,\ldots,U+B-1\}$, where $U\asymp N$ and $B=O(\log N)$.
Put
\[
 d_{N,B}=\left\lceil\frac{\log U}{\log(64B)}\right\rceil.
\]
If the relative square-relation code has dimension $r\ge t$, then
\begin{equation}
 \pi(B)+1\ge
 \sum_{i=0}^{t-1}\left\lceil d_{N,B}/2^i\right\rceil-t.
 \label{p1:limits:griesmer}
\end{equation}
Consequently $r\le1$ when $B=\log N/\log2+o(\log N)$.
More generally, if $B\sim c\log N$ and
$c<2-2^{1-t}$, then $r<t$ for all sufficiently large $N$.
For every fixed $\varepsilon>0$, the relative signs are exactly
independent in all windows with $B\le(1-\varepsilon)\log N$.
\end{lemma}
\begin{proof}
Every square relation is supported on the $D$ integers whose odd
prime valuations are all at primes at most $B$: a larger prime is
private in the window.  The even square-relation code on these
coordinates has redundancy at most $\pi(B)+1$.  Its minimum nonzero
weight is at least $d_{N,B}$ by Lemma~\ref{p1:arith:runge}.

For completeness, a binary linear code of length $D$, dimension $r$
and minimum distance $d$ has redundancy at least
$\sum_{i<t}\lceil d/2^i\rceil-t$ whenever $r\ge t$.
Choose a word of minimum weight $d_0\ge d$ and project outside its
support.  The kernel of this projection is its one-dimensional span:
another nonzero word supported there would have smaller weight.
The projected code has dimension $r-1$ and minimum distance at least
$\lceil d_0/2\rceil$, since one of two complementary lifts uses at
most $\lfloor d_0/2\rfloor$ coordinates on that support.  The
redundancy decreases by $d_0-1$.  Repeating $t$ times proves the
claim, and hence \eqref{p1:limits:griesmer}.

The prime number theorem gives
$\pi(B)\sim c\log N/\log\log N$, whereas
$d_{N,B}\sim\log N/\log\log N$ when $B\sim c\log N$.
Thus \eqref{p1:limits:griesmer} excludes dimension $t$ under the
stated strict inequality.  At critical length, $1/\log2<3/2$,
so $t=2$ applies.  The assertion for every smaller subcritical
window follows by including it in a window of length
$\lfloor(1-\varepsilon/2)\log N\rfloor$ and projecting.
\end{proof}

